\documentclass[12pt,a4paper]{article}

\usepackage[utf8]{inputenc}
\usepackage[T1]{fontenc}
\usepackage[a4paper, top=3cm, bottom=2cm, left=3cm, right=2cm]{geometry}
\usepackage{setspace}
\usepackage{indentfirst}
\usepackage{times}
\usepackage{csquotes}
\usepackage[hidelinks]{hyperref}
\usepackage{titlesec}
\usepackage[round,authoryear]{natbib}

\hyphenation{mu-si-co-te-ra-pia mu-si-co-te-ra-pêu-tica fe-no-me-no-ló-gi-ca
me-ta-ma-te-rial e-le-tro-mag-né-ti-cas com-po-si-tor com-po-si-cio-nal
in-ven-ção in-ven-ções pe-da-go-gi-ca pe-da-gó-gi-co tem-po-ra-li-da-de
mu-si-cal me-lódi-ca har-mô-ni-ca contraponto contrapontístico
es-tro-fi-ca-ção or-ni-to-lo-gi-ca trans-cri-ção transcrições polifo-nia
inter-pre-ta-ti-va paralelismo Hervieu-Léger Schopenhauer Jankélévitch
Wahrheitsgehalt scrutoniano jankélévitchiano paralelismos sinonímico
antitético sintético acústicas musicoterapêutica musicoterapêutico
Benenzon Bruscia Schneider pitagóricos boecianos boeciana}

\onehalfspacing
\setlength{\parindent}{1.25cm}

\titleformat{\section}{\normalfont\bfseries\large}{\thesection}{1em}{}
\titleformat{\subsection}{\normalfont\bfseries}{\thesubsection}{1em}{}
\titlespacing*{\section}{0pt}{18pt}{12pt}
\titlespacing*{\subsection}{0pt}{14pt}{8pt}

\hypersetup{
  pdftitle={A figura de Davi na tradicao musicoterapeutica ocidental},
  pdfauthor={Eduardo de Souza},
  pdfsubject={Musicoterapia, filosofia da musica e educacao musical}
}

\begin{document}

\begin{center}
{\Large\bfseries A figura de Davi na tradição musicoterapêutica ocidental:}\\[0.5em]
{\large leitura em três ciclos articulada com a filosofia da música contemporânea\\e aplicação pedagógica na educação básica brasileira}\\[1.5em]
{Eduardo de Souza\footnote{Licenciado em Educação Artística com habilitação em Música pela Faculdade Mozarteum de São Paulo (FAMOSP, 2024). Mestre em Engenharia Elétrica pela Pontifícia Universidade Católica de Minas Gerais (PUC Minas, 2018). Professor da rede municipal de Contagem, Minas Gerais. Contato: \texttt{edueandreia@gmail.com}.}}\\[0.5em]
{\small Maio de 2026}
\end{center}

\vspace{1em}

\noindent\textbf{Resumo}\\[0.5em]
{\small\noindent O presente artigo investiga a figura bíblica de Davi como paradigma fundador da tradição musicoterapêutica ocidental, propondo leitura sistemática em três ciclos articulados: o ciclo da escuta (Davi pastor das colinas de Belém), o ciclo da invenção (Davi salmista, autor formal dos Salmos canônicos) e o ciclo do cuidado (Davi tocando harpa para o rei Saul atormentado, conforme 1 Samuel 16). Os três ciclos são analisados em sua continuidade interna e em sua reatualização possível na educação musical básica brasileira contemporânea. A reflexão articula-se com cinco vozes da filosofia da música ocidental --- Schopenhauer, Jankélévitch, Adorno, Scruton e Wisnik --- e com a tradição musicoterapêutica institucionalizada por Bruscia, Benenzon e Schneider. O artigo propõe que a figura davídica, longe de constituir apenas referência simbólica retórica nos manuais clínicos da disciplina, oferece estrutura analítica operacional para pensar a relação entre escuta-mundo, invenção-forma e cuidado-comunidade na prática musical contemporânea.}

\vspace{0.5em}
\noindent\textbf{Palavras-chave:} Davi; harpa; salmos; musicoterapia; paisagem sonora; educação musical; filosofia da música; cuidado.

\vspace{1.5em}

\begin{flushright}
\begin{minipage}{0.65\textwidth}
\itshape\small
``Se a Arte não vem de dentro, não é arte: é expressão de circunstância.''

\upshape\rmfamily\normalsize\flushright
--- Ronaldo Cupertino da Silva, artista plástico\\
(comunicação pessoal ao autor, com autorização formal)
\end{minipage}
\end{flushright}

\vspace{1em}

\section{Introdução: por que retornar a Davi}

A história da musicoterapia ocidental, em sua versão narrada pelos manuais introdutórios da disciplina, costuma abrir com a passagem de 1 Samuel 16, na qual o jovem Davi é convocado à corte do rei Saul para acalmar, por meio da harpa, o espírito perturbado do monarca. A passagem é citada com tal frequência que se tornou, na prática, lugar-comum disciplinar: \citet{Bruscia2000} a menciona em \emph{Definindo Musicoterapia}, e a maioria dos manuais introdutórios da disciplina recorre à mesma referência. A presença é unânime; o tratamento, contudo, é raso. Davi aparece como anedota fundadora, exemplo retórico para introduzir a disciplina ao leigo, raramente como objeto de análise sistemática que pudesse iluminar a estrutura conceitual da própria musicoterapia.

O presente artigo propõe-se a corrigir essa lacuna em pequena escala, oferecendo leitura da figura davídica em três ciclos articulados (escuta, invenção, cuidado) e demonstrando que esses três ciclos, longe de constituírem distribuição arbitrária de momentos biográficos, organizam de modo coerente as três dimensões fundamentais da prática musical em registro terapêutico, pedagógico e artístico. A hipótese é que Davi não é apenas exemplo histórico, mas estrutura analítica.

O retorno à figura davídica justifica-se por três razões. A primeira é histórica: trata-se da referência mais antiga e contínua da tradição musicoterapêutica ocidental. A segunda é estrutural: a articulação entre escuta-mundo, invenção-forma e cuidado-comunidade que se observa na narrativa de Davi reaparece, sob diferentes vocabulários, em todas as principais escolas musicoterapêuticas contemporâneas, sugerindo que se trata de estrutura profunda da própria prática. A terceira é pedagógica: a reatualização desses três ciclos em escala adequada à educação musical básica oferece caminho concreto para a formação de professores que articulem dimensão artística e dimensão de cuidado em sua prática docente.

\section{Revisão bibliográfica}

Este trabalho mobiliza quatro corpos literários distintos. A revisão aqui apresentada é seletiva, focando os autores efetivamente operados ao longo do artigo.

\subsection{Tradição musicoterapêutica}

O campo da musicoterapia foi sistematizado academicamente na segunda metade do século XX. \citet{Bruscia2000}, em \emph{Definindo Musicoterapia}, oferece síntese institucional do campo, com proposta de definição operacional que se tornou referência internacional. Rolando Benenzon, em sua produção argentina, desenvolveu o conceito de \emph{identidade sonora} (ISO) que estrutura abordagem clínica largamente adotada na América Latina.

\subsection{Estudos bíblicos e exegese de 1 Samuel}

Para os propósitos deste artigo, mobiliza-se principalmente \citet{Biblia1993} como texto-fonte, com atenção especial a 1 Samuel 16 (passagem de Davi e Saul) e ao Livro dos Salmos. A exegese histórico-crítica do material é apoiada pelas introduções clássicas aos estudos do Antigo Testamento, com atenção à datação provável do Reino unificado (séc. X a.C.).

\subsection{Paisagem sonora e ecomusicologia}

O conceito de \emph{paisagem sonora} (\emph{soundscape}), desenvolvido por \citet{Schafer1994, Schafer1991}, fornece o aparato analítico para pensar a relação entre Davi pastor e seu ambiente acústico. A noção de \emph{acustemologia} de \citet{Feld1982} oferece base antropológica complementar. No contexto brasileiro, \citet{Fonterrada2008} articula a tradição schaferiana com a pedagogia musical nacional, e \citet{Wisnik1989} oferece a articulação som-pulso-altura que estrutura a análise musical aqui empreendida.

\subsection{Filosofia da música ocidental contemporânea}

Cinco autores fornecem o quadro filosófico: \citet{Schopenhauer2005}, com a tese da música como expressão direta da Vontade; \citet{Jankelevitch1983}, com a noção de inefável musical; \citet{Adorno2018}, com o conceito de conteúdo de verdade (\emph{Wahrheitsgehalt}); \citet{Scruton1997}, com a tese da intencionalidade musical; e o já citado Wisnik. A esses se soma John \citet{Cage1961}, cuja meditação sobre silêncio informa o tratamento do tempo musical aqui proposto.

\section{Davi: contexto histórico e exegese do nome}

\subsection{Contexto histórico do Reino unificado}

Conforme reconstituição predominante na historiografia bíblica contemporânea, Davi teria reinado em Judá e depois sobre o Israel unificado em torno do século X a.C., sucedendo Saul e sendo sucedido por Salomão. O período é caracterizado pela transição entre a estrutura tribal das doze tribos de Israel e o estabelecimento de uma monarquia centralizada com capital em Jerusalém. A figura de Davi, nessa transição, é simultaneamente militar, política e cultual --- guerreiro, rei e organizador da liturgia musical do templo.

\subsection{Exegese do nome}

O nome hebraico Davi é etimologicamente conectado à raiz \emph{dwd}, que carrega o sentido nuclear de ``amar'', ``ser amado'' ou ``amado''. A literatura exegética discute se o nome significa ``amado'' (forma passiva) ou ``amante'' (forma ativa). A maior parte dos comentadores favorece a primeira interpretação. A configuração afetiva inscrita no próprio nome do personagem ressoa com a dimensão de cuidado que se observa em sua atuação musical.

\subsection{Davi como personagem literário multifacetado}

O Davi dos textos canônicos é personagem complexo, longe da idealização pia que algumas leituras devocionais sugerem. Os mesmos textos que registram sua capacidade musical e seu favor divino registram também adultério, homicídio indireto, conflitos familiares severos e oscilações de fé. Para os fins desta análise, será privilegiado o Davi de 1 Samuel 16 a 2 Samuel 5 --- o jovem pastor, o salmista, o músico que toca para Saul.

\section{Primeiro ciclo --- A escuta: Davi pastor}

\subsection{O lugar do pastoreio}

Antes do palácio, antes da composição, antes do conflito com Saul, Davi é pastor. A tradição o apresenta menino no campo, sozinho com o rebanho do pai Jessé nas colinas de Belém. A formação musical de Davi não começa em academia nem em mestre humano. Começa em horas e horas de exposição diária a uma paisagem sonora particular: o balido das ovelhas, o vento entre as oliveiras, o canto dos pássaros das colinas da Judeia, o silêncio espesso das noites do deserto.

\subsection{A paisagem sonora das colinas de Belém}

As colinas em torno de Belém, situadas a cerca de oito quilômetros ao sul de Jerusalém, apresentam topografia íngreme, vegetação esparsa, presença significativa de oliveiras. O som dominante de uma jornada típica de pastoreio nessas colinas seria, provavelmente, o balido constante das ovelhas, sobreposto ao som do vento, com aves canoras pontuando o espectro médio-agudo, e atravessando tudo, períodos prolongados de silêncio relativo.

Essa paisagem sonora, sustentada por horas, dias, anos, constitui o que Schafer denominaria \emph{keynote sound} --- o som-tônica de fundo que estrutura toda a percepção auditiva do sujeito que nele habita \citep{Schafer1994}.

\subsection{O pastor como ouvinte privilegiado}

A condição de pastor confere a Davi posição auditiva privilegiada por três razões. Primeira, isolamento: longos períodos sozinho com o rebanho permitem grau de atenção auditiva que vida urbana raramente proporciona. Segunda, vigília: o pastoreio exige escuta ativa para detectar predadores, animais perdidos, mudanças climáticas. Terceira, repetição: anos de exposição ao mesmo tipo de paisagem sonora permitem desenvolver familiaridade íntima com seus padrões.

\subsection{A formação musical sem mestre}

Um ponto que merece destaque: os textos canônicos não registram nenhum mestre musical humano para Davi. A formação musical aparece como algo que Davi já possui quando entra em cena. O que isso sugere é que a formação musical de Davi se deu primariamente por imersão prolongada na paisagem sonora natural do pastoreio.

Como educador musical e estudante das relações entre música e cuidado humano, considero que a harpa de Davi foi a expressão máxima de algo que veio de dentro para expressar o sentimento natural --- perfeito --- do Poeta Supremo, que é Deus. Essa expressão máxima começou, primeiro, em escuta.

\section{Segundo ciclo --- A invenção: Davi salmista}

\subsection{Os Salmos como invenções}

Da escuta do pastor brota, em algum momento da biografia davídica, o salmista. A tradição atribui a Davi cerca de setenta e três dos cento e cinquenta Salmos canônicos. Os Salmos não são improvisos despreocupados. São construções formais cuidadosas, com paralelismos hebraicos rigorosos, refrões estruturais, estruturas estróficas regulares.

\subsection{O paralelismo hebraico como estrutura compositiva}

A poesia hebraica antiga opera primariamente por paralelismo --- repetição com variação de uma mesma ideia em duas linhas sucessivas. O paralelismo sinonímico repete a mesma ideia em palavras diferentes (Salmo 23:1-2). O paralelismo antitético contrasta duas ideias opostas (Salmo 1:6). O paralelismo sintético desenvolve, expande ou complementa a primeira linha na segunda (Salmo 93:1).

O efeito musical do paralelismo é considerável: na execução cantada, ele oferece estrutura natural para frases musicais simétricas, com perguntas e respostas. O paralelismo é, em rigor, ferramenta de cuidado pedagógico --- facilita memorização, internalização, transmissão.

\subsection{A invenção que responde, não cria do nada}

A invenção de Davi tem uma característica que a distingue: ela não inventa do nada, ela responde. Cada salmo é resposta a algo previamente escutado --- a misericórdia divina, a perseguição inimiga, o lamento próprio, a glória da criação.

De minha perspectiva pessoal, formada na dupla experiência da música como prática herdada (recebida do meu pai Izaias Maciel, músico em atividade) e como objeto de reflexão acadêmica, a harpa de Davi parece a expressão máxima de algo que vem de dentro para encontrar o sentimento natural --- perfeito --- do Poeta Supremo, que é Deus.

\subsection{Tipologia dos Salmos}

A tradição crítica costuma classificar os Salmos em algumas categorias principais: lamentos individuais, lamentos coletivos, hinos, salmos de ações de graças, salmos reais, salmos de Sião, salmos de sabedoria. Essa tipologia revela algo importante: a prática davídica de invenção não se restringe a um único registro afetivo. Davi articula a gama afetiva inteira da experiência humana.

\section{Terceiro ciclo --- O cuidado: Davi e Saul}

\subsection{A passagem fundadora de 1 Samuel 16}

A passagem que mais imediatamente associa Davi à musicoterapia é 1 Samuel 16:14-23 \citep{Biblia1993}. O texto descreve a condição de Saul como ``espírito maligno proveniente de Deus'' --- formulação teologicamente complexa que tem suscitado debates por dois milênios.

Importante: o texto não sugere que Davi cure Saul de modo definitivo. A formulação é precisa: quando Davi toca, Saul ``sentia alívio'', ``se achava melhor'', e o espírito maligno ``se retirava dele''. A musicoterapia davídica opera por alívio sintomático recorrente, não por cura definitiva.

\subsection{A prática: ``tomava Davi a harpa e a dedilhava''}

A descrição da prática musical de Davi para Saul é sóbria, quase telegráfica: tomava a harpa e a dedilhava. Não há descrição do repertório, da duração das sessões, da postura, do ambiente. O texto retém o essencial: Davi tocava, e Saul melhorava. Essa sobriedade descritiva sugere que o ato musical fundamental --- alguém que toca para alguém que sofre --- opera por dinâmicas que excedem qualquer roteiro específico.

\subsection{O que cura: não a habilidade, mas a profundidade}

O que torna a música de Davi capaz desse efeito? Não é virtuosismo técnico isolado. O que torna a música de Davi capaz de aliviar Saul é a continuidade com os dois ciclos anteriores: ela é música de alguém que primeiro escutou o mundo, depois aprendeu a responder em forma. O que alivia Saul não é a habilidade de Davi --- é a profundidade do que Davi traz.

\subsection{A relação Davi-Saul como modelo do encontro terapêutico}

Davi é jovem, Saul é o rei estabelecido. Há, na superfície, assimetria de poder maciça em favor de Saul. Mas no plano específico do encontro musical-terapêutico, a configuração se inverte: Davi é o cuidador, Saul é o cuidado. Essa inversão temporária e localizada de hierarquia é constitutiva de toda relação terapêutica genuína.

\section{Síntese: os três ciclos como movimento único}

A tese central deste artigo é que os três ciclos --- escuta, invenção, cuidado --- não constituem etapas sucessivas no tempo, mas dimensões simultâneas e mutuamente sustentadas de uma única prática musical integral.

Davi não é primeiro pastor, depois salmista, depois cuidador, como se cada fase substituísse a anterior. Ele é, em momentos diversos de sua vida e até simultaneamente, as três coisas. A estrutura é espiralada: cada ciclo sustenta o próximo e é retroalimentado por ele. A escuta sustenta a invenção; a invenção sustenta o cuidado; o cuidado retroalimenta a escuta.

\section{Recepção de Davi na tradição musicoterapêutica}

A presença de Davi nos manuais introdutórios da musicoterapia ocidental é praticamente universal. \citet{Bruscia2000} menciona o episódio de 1 Samuel 16 como exemplo histórico precoce de aplicação intencional da música para fins terapêuticos. Essa unanimidade contrasta, contudo, com a frequente superficialidade de seu tratamento.

No contexto brasileiro, a musicoterapia profissionalizou-se ao longo da segunda metade do século XX, com regulamentação profissional pela Lei nº 14.842 de 11 de abril de 2024. A relevância contemporânea de Davi para a musicoterapia brasileira, neste artigo, é menos a de modelo clínico operacional e mais a de estrutura conceitual que pode iluminar tanto a prática clínica quanto suas extensões pedagógicas em educação musical básica.

\section{Diálogo com a filosofia da música contemporânea}

\subsection{Schopenhauer e a Vontade}

\citet{Schopenhauer2005} atribui à música estatuto único entre as artes: ela expressa diretamente a Vontade metafísica, sem mediação de conceito. Aplicada à harpa de Davi, a tese schopenhaueriana sugere que o efeito de alívio sobre Saul não é mediado por conteúdo conceitual, mas opera diretamente sobre a estrutura volitiva do rei.

\subsection{Jankélévitch e o inefável}

\citet{Jankelevitch1983} recusa qualquer redução da música ao conceito. O texto não registra que Davi tenha pronunciado palavra alguma para Saul. A intervenção é puramente instrumental: Davi tomava a harpa e a dedilhava. Não há sermão, não há explicação. E justamente nessa abstinência da palavra opera o efeito.

\subsection{Adorno e o conteúdo de verdade}

\citet{Adorno2018} sustenta que a música possui \emph{Wahrheitsgehalt} --- conteúdo de verdade --- que se realiza no enfrentamento histórico entre forma e material. A harpa de Davi para Saul não opera por idealização ou maquiagem da situação do rei. Ela opera por enfrentamento da condição real.

\subsection{Scruton e a intencionalidade musical}

\citet{Scruton1997} defende que ouvir música é ato intencional distinto de ouvir mero som. A música de Davi não muda os sons que cercam Saul, mas modifica a intencionalidade com que Saul escuta. O alívio que Saul sente não é mudança do mundo em volta --- é mudança do seu modo de escutar o mundo.

\subsection{Wisnik e a tridimensionalidade do som}

\citet{Wisnik1989} articula as três dimensões físicas do som --- ruído, pulso, altura --- com as três dimensões formais da música. A harpa de Davi opera com equilíbrio particular entre as três dimensões, sem privilegiar excessivamente a altura, preservando no som davídico o ruído originário do mundo.

\subsection{Síntese filosófica}

Articulando essas cinco vozes, a tese central pode ser enunciada em formulação técnica: a harpa de Davi opera como expressão direta de uma interioridade musical à maneira de Schopenhauer; no registro do inefável jankélévitchiano; sob a exigência adorniana de honestidade histórica; modificando intencionalmente a escuta no sentido scrutoniano; e preservando a tridimensionalidade som-pulso-altura descrita por Wisnik.

\section{Aplicação na educação musical básica brasileira}

A leitura da figura davídica em três ciclos tem aplicação direta em prática pedagógica na educação musical básica, conforme o autor tem desenvolvido em sua atuação como professor concursado da rede municipal de Contagem, Minas Gerais.

\subsection{Cartografia sonora da escola (ciclo da escuta)}

O primeiro exercício consiste em propor aos estudantes que percorram os espaços da escola com a única tarefa de escutar, sem julgar nem nomear. Ao retornarem à sala, os estudantes registram, em notação livre, três sons que mais lhes chamaram a atenção. O conjunto desses registros constitui uma cartografia sonora coletiva da escola.

Esse exercício é a reatualização, em escala adequada à educação básica, do primeiro ciclo davídico: escutar o mundo antes de tocá-lo.

\subsection{Invenção sobre motivo encontrado (ciclo da invenção)}

O segundo exercício parte da cartografia sonora. Cada estudante escolhe um som específico e o transcreve melodicamente. Uma vez transcrito, o motivo é trabalhado por procedimentos elementares de variação: repetição, deslocamento de altura, inversão, expansão rítmica. O resultado final é uma pequena invenção autoral.

Quando o estudante transcreve um som escutado e o desdobra em pequena invenção, ele refaz o gesto do salmista em escala miniaturizada.

\subsection{Coro de paisagem (ciclo do cuidado)}

O terceiro exercício é o mais avançado. Em formação coletiva, grupos pequenos reproduzem vocalmente elementos da paisagem (vento, chuva, pássaro), organizados em sequência ou sobreposição. O aprendizado mais importante do exercício é relacional: cada voz precisa se ajustar às outras, cada gesto sonoro precisa responder aos demais. Os estudantes experimentam a função do cuidado mútuo via música.

\subsection{Observações sobre a prática}

A aplicação dos três exercícios gerou observações que merecem registro: estudantes inicialmente identificados como ``sem talento musical'' apresentaram performances notáveis nos exercícios de cartografia sonora. A sequência três-ciclos produziu queda significativa em conflitos interpessoais durante as aulas. A articulação explícita com a figura de Davi (apresentada como histórico-cultural, sem proselitismo religioso) ofereceu aos estudantes ancoragem narrativa que facilitou compreensão das três fases do trabalho.

\section{Considerações finais}

Este artigo argumentou que a figura de Davi oferece estrutura analítica de três ciclos articulados --- escuta, invenção, cuidado --- capaz de iluminar tanto a prática clínica musicoterapêutica regulamentada quanto suas extensões pedagógicas em educação musical básica. A leitura dos três ciclos como movimento único e simultâneo permitiu mostrar a continuidade interna do paradigma davídico, em diálogo articulado com cinco vozes da filosofia da música.

A aplicação pedagógica em três exercícios testados em sala de aula na rede municipal de Contagem demonstrou que a estrutura analítica é operacionalizável em escala adequada à educação básica brasileira. Essa operacionalização não pretende ser musicoterapia clínica --- atividade regulamentada pela Lei nº 14.842/2024 e que exige formação universitária específica --- mas educação musical informada pela tradição musicoterapêutica.

Fecho com a frase de meu irmão Ronaldo Cupertino da Silva, citada nesta obra com sua autorização formal:

\begin{quote}
\itshape Se a Arte não vem de dentro, não é arte: é expressão de circunstância.
\end{quote}

Em todos os três ciclos, a circunstância é ocasião, mas não fonte. A fonte está dentro, e a tarefa do músico, do educador e do cuidador, em qualquer época, é a mesma: oferecer ao outro aquilo que vem de dentro.

\renewcommand{\refname}{Referências}
\begin{thebibliography}{99}

\bibitem[Adorno(2018)]{Adorno2018}
ADORNO, T. W. \textit{Filosofia da nova música}. Trad. M. Bandeira. São Paulo: Perspectiva, 2018.

\bibitem[Bíblia(1993)]{Biblia1993}
BÍBLIA SAGRADA. \textit{Almeida Revista e Atualizada}. Barueri: Sociedade Bíblica do Brasil, 1993. Particularmente: 1 Samuel 16; Livro dos Salmos.

\bibitem[Bruscia(2000)]{Bruscia2000}
BRUSCIA, K. E. \textit{Definindo musicoterapia}. Trad. M. Ronchini. Rio de Janeiro: Enelivros, 2000.

\bibitem[Cage(1961)]{Cage1961}
CAGE, J. \textit{Silence: Lectures and Writings}. Middletown: Wesleyan University Press, 1961.

\bibitem[Feld(1982)]{Feld1982}
FELD, S. \textit{Sound and Sentiment: Birds, Weeping, Poetics, and Song in Kaluli Expression}. Philadelphia: University of Pennsylvania Press, 1982.

\bibitem[Fonterrada(2008)]{Fonterrada2008}
FONTERRADA, M. T. de O. \textit{De tramas e fios: um ensaio sobre música e educação}. São Paulo: Editora UNESP, 2008.

\bibitem[Jankélévitch(1983)]{Jankelevitch1983}
JANKÉLÉVITCH, V. \textit{La musique et l'ineffable}. Paris: Seuil, 1983.

\bibitem[Schafer(1991)]{Schafer1991}
SCHAFER, R. M. \textit{O ouvido pensante}. Trad. M. T. Fonterrada et al. São Paulo: Editora UNESP, 1991.

\bibitem[Schafer(1994)]{Schafer1994}
SCHAFER, R. M. \textit{The Soundscape: Our Sonic Environment and the Tuning of the World}. Rochester: Destiny Books, 1994.

\bibitem[Schopenhauer(2005)]{Schopenhauer2005}
SCHOPENHAUER, A. \textit{O mundo como vontade e representação}. Trad. J. Barboza. São Paulo: Unesp, 2005.

\bibitem[Scruton(1997)]{Scruton1997}
SCRUTON, R. \textit{The Aesthetics of Music}. Oxford: Oxford University Press, 1997.

\bibitem[Silva(s.d.)]{Silva}
SILVA, R. C. da. Comunicação pessoal ao autor. [s.d.]. Citada nesta obra mediante autorização formal do citado.

\bibitem[Souza(2018)]{Souza2018}
SOUZA, E. de. \textit{Projeto Ótimo de uma Antena Patch com Metamaterial}. 2018. 79 f. Dissertação (Mestrado em Engenharia Elétrica) --- Pontifícia Universidade Católica de Minas Gerais, Belo Horizonte, 2018. Orientadora: Rose Mary de Souza Batalha. Disponível em: \url{https://bib.pucminas.br/teses/EngEletrica_SouzaE_1.pdf}.

\bibitem[Wisnik(1989)]{Wisnik1989}
WISNIK, J. M. \textit{O som e o sentido: uma outra história das músicas}. São Paulo: Companhia das Letras, 1989.

\end{thebibliography}

\end{document}
